\begin{table}[!ht]
\centering
\caption{Decomposição de Oaxaca--Blinder (twofold) do gap salarial racial --- especificação
de \emph{acesso} (com ocupação e contexto de UPA). Efeito dotação: parcela explicada por
diferenças em características observáveis, \emph{incluindo a posição ocupacional}. Efeito
coeficiente: parcela não explicada (limite inferior da discriminação salarial \emph{dentro}
da ocupação). População completa da PEA; erros-padrão por bootstrap de 200 replicações
agrupadas por UF. Ressalva (Oaxaca \& Ransom, 1999): tratar a ocupação como dotação tende a
\emph{subestimar} a discriminação total, pois a segregação ocupacional é, ela própria,
discriminatória --- daí a complementaridade com o GLMM de acesso (Tabela~\ref{tab:glmm_glassceil}).}
\label{tab:oaxaca_blinder}
\begin{tabular}{llcc}
\toprule
Componente & Magnitude & \% do Gap & SE (bootstrap) \\
\midrule
Gap Total (branco $-$ negro) & 0,4255 & 100,0 & --- \\
Efeito Dotação (características, incl.\ ocupação) & 0,3566 & 83,8 & (0,0288) \\
Efeito Coeficiente (não explicado / discriminação) & 0,0689 & 16,2 & (0,0121) \\
\bottomrule
\end{tabular}
\end{table}
